% !TeX program = pdflatex
% !TeX encoding = UTF-8
\documentclass[11pt,aspectratio=169,t]{beamer}
% Все презентации курса компилируются только pdflatex.
\usepackage{iftex}
\RequirePDFTeX
\usepackage{cmap}
\usepackage[T2A]{fontenc}
\usepackage[utf8]{inputenc}
\usepackage[english,russian]{babel}
\usepackage{PTSerif}
\usepackage{amsmath,amsfonts,amssymb}
\usepackage{graphicx,microtype,anyfontsize,tikz}
\usepackage{ragged2e}
\usetikzlibrary{arrows.meta,calc}

\usetheme{default}
\usefonttheme{serif}
\DeclareMathOperator{\sen}{sen}
\bibliographystyle{apalike}

% Общие сведения; номер лекции, её название и дата задаются в N.tex.
\newcommand{\coursename}{Введение в маневрирование КА}
\author{Чиняев Владимир Николаевич}
\institute{Московский физико-технический институт}

% Книжное оформление с голубым акцентом.
\definecolor{coursepaper}{HTML}{FBF9F3}
\definecolor{courseink}{HTML}{302B28}
\definecolor{courseaccent}{HTML}{187AA5}
\definecolor{coursemuted}{HTML}{6D6560}
\definecolor{courserule}{HTML}{A9CADA}
\definecolor{coursetint}{HTML}{EDF4F6}
\tikzset{
  course diagram/.style={font=\fontsize{9.5}{11.5}\selectfont,
    line cap=round,line join=round,>=Latex},
  orbit line/.style={draw=courseaccent,line width=1pt},
  aux line/.style={draw=courserule,line width=.5pt},
  vector line/.style={draw=courseink,line width=.65pt,->},
  angle line/.style={draw=courseaccent,line width=.8pt,->},
  orbit point/.style={circle,fill=courseaccent,inner sep=1.7pt}
}
\setbeamercolor{normal text}{fg=courseink,bg=coursepaper}
\setbeamercolor{background canvas}{bg=coursepaper}
\setbeamercolor{structure}{fg=courseaccent}
\setbeamercolor{alerted text}{fg=courseaccent}
\setbeamercolor{frametitle}{fg=courseink,bg=coursepaper}
\setbeamercolor{footline}{fg=coursemuted,bg=coursepaper}
\setbeamercolor{caption name}{fg=courseaccent}

% Основной текст — 11/14 pt; подписи к рисункам — 9.5/11.5 pt.
\let\coursenormalsize\normalsize
\renewcommand{\normalsize}{\coursenormalsize\fontsize{11}{14}\selectfont}
\setbeamerfont{normal text}{size*={11pt}{14pt}}
\setbeamerfont{frametitle}{size*={20pt}{23pt},series=\mdseries}
\setbeamerfont{title}{size*={27pt}{31pt},series=\mdseries}
\setbeamerfont{subtitle}{size*={12.8pt}{15pt},series=\mdseries}
\setbeamerfont{author}{size*={10.8pt}{13pt}}
\setbeamerfont{institute}{size*={9.2pt}{11pt}}
\setbeamerfont{date}{size*={9pt}{11pt}}
\setbeamerfont{footline}{size*={6.8pt}{8pt}}
\setbeamerfont{caption}{size*={9.5pt}{11.5pt}}
\setbeamerfont{itemize/enumerate body}{size*={11pt}{14pt}}
\setbeamerfont{itemize/enumerate subbody}{size*={11pt}{14pt}}
\setbeamerfont{itemize/enumerate subsubbody}{size*={11pt}{14pt}}
\AtBeginDocument{\normalsize}

\setbeamersize{text margin left=6mm,text margin right=6mm}
\setbeamertemplate{navigation symbols}{}
\setbeamertemplate{headline}{}
\setbeamertemplate{caption}[numbered]
\setbeamertemplate{itemize items}[circle]

% Стандартные блоки Beamer позволяют группировать текст и формулы.
\setbeamertemplate{blocks}[default]
\setbeamerfont{block title}{size*={11pt}{14pt},series=\bfseries}
\setbeamerfont{block body}{size*={11pt}{14pt}}
\setbeamercolor{block title}{fg=courseaccent,bg=coursepaper}
\setbeamercolor{block body}{fg=courseink,bg=coursetint}
\setbeamercolor{block title alerted}{parent=block title}
\setbeamercolor{block body alerted}{parent=block body}
\setbeamercolor{block title example}{parent=block title}
\setbeamercolor{block body example}{parent=block body}

% Компактный заголовок; его высота учитывается при размещении содержания.
\setbeamertemplate{frametitle}{%
  \vskip4.5mm%
  \begin{beamercolorbox}[wd=\textwidth,sep=0pt]{frametitle}
    \usebeamerfont{frametitle}\strut\insertframetitle\strut\par
    \vskip2mm%
    {\color{courseaccent}\hrule height .55pt}
  \end{beamercolorbox}%
}

% Нижний колонтитул: курс слева, лекция и номер слайда справа.
\setbeamertemplate{footline}{%
  \begin{beamercolorbox}[wd=\paperwidth,ht=5.5mm,dp=0pt]{footline}
    \begin{tikzpicture}[remember picture,overlay,x=1mm,y=1mm]
      \begin{scope}[shift={(current page.south west)}]
        \draw[courserule,line width=.3pt] (6,5.5)--(154,5.5);
        \node[text=coursemuted,anchor=west,inner sep=0pt,
          font=\usebeamerfont{footline}] at (6,2.8) {\coursename};
        \node[text=coursemuted,anchor=east,inner sep=0pt,
          font=\usebeamerfont{footline}] at (154,2.8)
          {Лекция \arabic{lecture}\quad\textbullet\quad\insertframenumber};
      \end{scope}
    \end{tikzpicture}%
  \end{beamercolorbox}%
}

% Титульник рассчитан на 16:9 (160 × 90 мм), без нижнего колонтитула.
\setbeamertemplate{title page}{%
  \begin{tikzpicture}[remember picture,overlay,x=1mm,y=1mm]
    \begin{scope}[shift={(current page.south west)}]
      \node[text=coursemuted,anchor=north,font=\usebeamerfont{institute}]
        at (80,82) {\insertinstitute};
      \draw[courseaccent,line width=.55pt] (71,71)--(89,71);
      \node[text=courseink,align=center,text width=125mm,
        font=\usebeamerfont{title}] at (80,53)
        {\hyphenpenalty=10000\coursename\par};
      \node[text=courseaccent,font=\usebeamerfont{subtitle}]
        at (80,32) {Лекция \arabic{lecture}. \inserttitle};
      \node[text=courseink,font=\usebeamerfont{author}]
        at (80,19) {\insertauthor};
      \node[text=coursemuted,font=\usebeamerfont{date}]
        at (80,9) {\insertdate};
    \end{scope}
  \end{tikzpicture}%
}

\usepackage{booktabs,array,tabularx}

\setcounter{lecture}{2}
\title{Методы прогноза баллистического движения}
\date{15 сентября 2026 г.}

\begin{document}

% 1
\begin{frame}[plain]
\titlepage
\end{frame}

% 2
\begin{frame}{Вековые возмущения}
\centering% Составляющие возмущений и среднее изменение на двух интервалах.
\begin{tikzpicture}[course diagram,x=1mm,y=1mm]
\path[use as bounding box] (-5,-5) rectangle (142,34);
\draw[vector line] (0,0)--(137,0) node[right] {$t$};
\draw[vector line] (0,0)--(0,30) node[above] {$c$};
% Вековая составляющая, сумма с долгой периодикой и полное изменение.
\draw[coursemuted,line width=.65pt] (0,10)--(134,{10+.033*134});
\draw[courserule,line width=1.1pt] plot[domain=0:134,samples=220]
 (\x,{10+.033*\x+11.5*sin(3.3*\x)});
\draw[courseaccent,line width=.7pt] plot[domain=0:134,samples=700]
 (\x,{10+.033*\x+11.5*sin(3.3*\x)+1.2*sin(42*\x)});
\foreach \x/\lab in {8/1,25/2,53/3,68/4}{
 \draw[courseink,line width=.45pt] (\x,-1)--(\x,{12+.033*\x+11.5*sin(3.3*\x)});
 \node[below] at (\x,-1) {$t_{\lab}$};
}
\coordinate (A) at (8,{10+.033*8+11.5*sin(26.4)});
\coordinate (B) at (25,{10+.033*25+11.5*sin(82.5)});
\coordinate (C) at (53,{10+.033*53+11.5*sin(174.9)});
\coordinate (D) at (68,{10+.033*68+11.5*sin(224.4)});
\draw[courseink,line width=1.5pt] (A)--(B) (C)--(D);
% Выноски среднего изменения заканчиваются точно на выделенных отрезках.
\node[align=center,text width=29mm,inner sep=0pt] (m1) at (18,29)
 {Среднее\\изменение};
\draw[vector line] (m1.south)--($(A)!.5!(B)$);
\node[align=center,text width=26mm,inner sep=0pt] (m2) at (65,19)
 {Среднее\\изменение};
\draw[vector line] (m2.south)--($(C)!.5!(D)$);
\node[align=left,anchor=north west,text width=85mm,inner sep=0pt] (sum) at (49,33)
 {Короткопериодические +\\долгопериодические + вековые};
% Выноска к гребню короткопериодического колебания, выше гладкой кривой.
\draw[vector line] ($(sum.south west)+(2,0)$)--
 (45,{10+.033*45+11.5*sin(148.5)+1.2*sin(1890)});
\node[align=center,inner sep=0pt] (sec) at (123,10.5) {Вековые};
\draw[vector line] (sec.north)--(123,{10+.033*123});
\node[align=center,text width=36mm,inner sep=0pt] (lp) at (121,4.2)
 {Долгопериодические\\и вековые};
% В этой точке полная кривая расположена выше гладкой: выноска их не пересекает.
\draw[vector line] (lp.west) to[out=170,in=0]
 (96,{10+.033*96+11.5*sin(316.8)});
\end{tikzpicture}
\par
\begin{block}{Определение}
\justifying
Вековые возмущения --- изменения орбитальных элементов, монотонно накапливающиеся со временем. Могут быть линейными по времени или пропорциональными какой-то степени времени.
\end{block}
\end{frame}

% 3
\begin{frame}{Долгопериодические возмущения}
\centering% Составляющие возмущений и среднее изменение на двух интервалах.
\begin{tikzpicture}[course diagram,x=1mm,y=1mm]
\path[use as bounding box] (-5,-5) rectangle (142,34);
\draw[vector line] (0,0)--(137,0) node[right] {$t$};
\draw[vector line] (0,0)--(0,30) node[above] {$c$};
% Вековая составляющая, сумма с долгой периодикой и полное изменение.
\draw[coursemuted,line width=.65pt] (0,10)--(134,{10+.033*134});
\draw[courserule,line width=1.1pt] plot[domain=0:134,samples=220]
 (\x,{10+.033*\x+11.5*sin(3.3*\x)});
\draw[courseaccent,line width=.7pt] plot[domain=0:134,samples=700]
 (\x,{10+.033*\x+11.5*sin(3.3*\x)+1.2*sin(42*\x)});
\foreach \x/\lab in {8/1,25/2,53/3,68/4}{
 \draw[courseink,line width=.45pt] (\x,-1)--(\x,{12+.033*\x+11.5*sin(3.3*\x)});
 \node[below] at (\x,-1) {$t_{\lab}$};
}
\coordinate (A) at (8,{10+.033*8+11.5*sin(26.4)});
\coordinate (B) at (25,{10+.033*25+11.5*sin(82.5)});
\coordinate (C) at (53,{10+.033*53+11.5*sin(174.9)});
\coordinate (D) at (68,{10+.033*68+11.5*sin(224.4)});
\draw[courseink,line width=1.5pt] (A)--(B) (C)--(D);
% Выноски среднего изменения заканчиваются точно на выделенных отрезках.
\node[align=center,text width=29mm,inner sep=0pt] (m1) at (18,29)
 {Среднее\\изменение};
\draw[vector line] (m1.south)--($(A)!.5!(B)$);
\node[align=center,text width=26mm,inner sep=0pt] (m2) at (65,19)
 {Среднее\\изменение};
\draw[vector line] (m2.south)--($(C)!.5!(D)$);
\node[align=left,anchor=north west,text width=85mm,inner sep=0pt] (sum) at (49,33)
 {Короткопериодические +\\долгопериодические + вековые};
% Выноска к гребню короткопериодического колебания, выше гладкой кривой.
\draw[vector line] ($(sum.south west)+(2,0)$)--
 (45,{10+.033*45+11.5*sin(148.5)+1.2*sin(1890)});
\node[align=center,inner sep=0pt] (sec) at (123,10.5) {Вековые};
\draw[vector line] (sec.north)--(123,{10+.033*123});
\node[align=center,text width=36mm,inner sep=0pt] (lp) at (121,4.2)
 {Долгопериодические\\и вековые};
% В этой точке полная кривая расположена выше гладкой: выноска их не пересекает.
\draw[vector line] (lp.west) to[out=170,in=0]
 (96,{10+.033*96+11.5*sin(316.8)});
\end{tikzpicture}
\par
\begin{block}{Определение}
\justifying
Долгопериодические возмущения --- периодические изменения орбитальных элементов, период которых существенно больше орбитального периода спутника (обычно, на один или два порядка).
\end{block}
\end{frame}

% 4
\begin{frame}{Короткопериодические возмущения}
\centering% Составляющие возмущений и среднее изменение на двух интервалах.
\begin{tikzpicture}[course diagram,x=1mm,y=1mm]
\path[use as bounding box] (-5,-5) rectangle (142,34);
\draw[vector line] (0,0)--(137,0) node[right] {$t$};
\draw[vector line] (0,0)--(0,30) node[above] {$c$};
% Вековая составляющая, сумма с долгой периодикой и полное изменение.
\draw[coursemuted,line width=.65pt] (0,10)--(134,{10+.033*134});
\draw[courserule,line width=1.1pt] plot[domain=0:134,samples=220]
 (\x,{10+.033*\x+11.5*sin(3.3*\x)});
\draw[courseaccent,line width=.7pt] plot[domain=0:134,samples=700]
 (\x,{10+.033*\x+11.5*sin(3.3*\x)+1.2*sin(42*\x)});
\foreach \x/\lab in {8/1,25/2,53/3,68/4}{
 \draw[courseink,line width=.45pt] (\x,-1)--(\x,{12+.033*\x+11.5*sin(3.3*\x)});
 \node[below] at (\x,-1) {$t_{\lab}$};
}
\coordinate (A) at (8,{10+.033*8+11.5*sin(26.4)});
\coordinate (B) at (25,{10+.033*25+11.5*sin(82.5)});
\coordinate (C) at (53,{10+.033*53+11.5*sin(174.9)});
\coordinate (D) at (68,{10+.033*68+11.5*sin(224.4)});
\draw[courseink,line width=1.5pt] (A)--(B) (C)--(D);
% Выноски среднего изменения заканчиваются точно на выделенных отрезках.
\node[align=center,text width=29mm,inner sep=0pt] (m1) at (18,29)
 {Среднее\\изменение};
\draw[vector line] (m1.south)--($(A)!.5!(B)$);
\node[align=center,text width=26mm,inner sep=0pt] (m2) at (65,19)
 {Среднее\\изменение};
\draw[vector line] (m2.south)--($(C)!.5!(D)$);
\node[align=left,anchor=north west,text width=85mm,inner sep=0pt] (sum) at (49,33)
 {Короткопериодические +\\долгопериодические + вековые};
% Выноска к гребню короткопериодического колебания, выше гладкой кривой.
\draw[vector line] ($(sum.south west)+(2,0)$)--
 (45,{10+.033*45+11.5*sin(148.5)+1.2*sin(1890)});
\node[align=center,inner sep=0pt] (sec) at (123,10.5) {Вековые};
\draw[vector line] (sec.north)--(123,{10+.033*123});
\node[align=center,text width=36mm,inner sep=0pt] (lp) at (121,4.2)
 {Долгопериодические\\и вековые};
% В этой точке полная кривая расположена выше гладкой: выноска их не пересекает.
\draw[vector line] (lp.west) to[out=170,in=0]
 (96,{10+.033*96+11.5*sin(316.8)});
\end{tikzpicture}
\par
\begin{block}{Определение}
\justifying
Короткопериодические возмущения --- периодические изменения орбитальных элементов, период которых сравним с орбитальным периодом спутника или меньше.
\end{block}
\end{frame}

% 5
\begin{frame}{Возмущения и точность прогноза}
\begin{block}{Вековые изменения}
\justifying
Приводят к неограниченному накоплению ошибки прогноза. Являются главной причиной деградации аналитических теорий на больших интервалах прогноза.
\end{block}
\begin{block}{Долгопериодические изменения}
\justifying
Обычно связаны с изменением ориентации линии узлов и линии апсид. Характерный период определяется орбитой и действующими возмущениями.
\end{block}
\begin{block}{Короткопериодические изменения}
\justifying
Обусловлены зависимостью возмущающих воздействий от быстрой переменной. При интегрировании необходимо разрешать эти изменения, поэтому они диктуют величину шага.
\end{block}
\end{frame}

% 6
\begin{frame}{Быстрые и медленные переменные}
\justifying
Необходимо разделять орбитальные элементы на быстрые и медленные переменные в зависимости от их относительной скорости изменения.
\begin{block}{Быстрые переменные}
\justifying
Переменные, существенно изменяющиеся за один орбитальный период даже в отсутствие возмущений. Например: аномалии $M$, $E$, $\nu$ изменяются за виток на $2\pi$; декартовы координаты.
\end{block}
\begin{block}{Медленные переменные}
\justifying
Переменные, изменения которых за один виток малы и обусловлены возмущающими воздействиями. В невозмущённой задаче такие элементы будут постоянными. Например: кеплеровы элементы орбиты $a$, $e$, $i$, $\Omega$, $\omega$.
\end{block}
\end{frame}

% 7
\begin{frame}{Оскулирующая орбита}
\begin{columns}[T,onlytextwidth]
\begin{column}{.51\textwidth}
\begin{block}{Определение}
\justifying
Оскулирующая орбита --- орбита задачи двух тел, определяемая фактическими положением и скоростью КА на рассматриваемую эпоху $t_*$. По ней КА двигался бы при исчезновении возмущающих сил.
\end{block}
\justifying
Её элементы называются \textbf{оскулирующими}. Они зависят от времени и содержат вековые, долго- и короткопериодические изменения.

\end{column}
\begin{column}{.47\textwidth}\centering
% Локальная волнистая дуга и касательная кеплерова орбита.
% При u=0 совпадают положение и направление касательной.
\begin{tikzpicture}[course diagram,x=.88mm,y=.88mm]
\path[use as bounding box] (-37,-27) rectangle (30,29);
\draw[orbit line] plot[domain=0:360,samples=220]
 ({21*cos(\x)/(1+.35*cos(\x))},{21*sin(\x)/(1+.35*cos(\x))});
\coordinate (P) at ({21*cos(40)/(1+.35*cos(40))},{21*sin(40)/(1+.35*cos(40))});
% Касательная единичного направления (-sin nu,e+cos nu), нормаль к ней.
\begin{scope}[shift={(P)},x={(-.43912mm,.76208mm)},y={(.76208mm,.43912mm)}]
\draw[courseink,dashed,line width=.85pt]
 plot[domain=-29:16,samples=180] (\x,{.007*\x^2+2.3*(1-cos(16*\x))});
\draw[vector line] (0,0)--(10,0) node[above left] {$\mathbf v(t_*)$};
\end{scope}
\filldraw[fill=coursetint,draw=courserule] (0,0) circle (2.6);
\node[below] at (0,-2.6) {Земля};
\draw[vector line] (0,0)--(P) node[pos=.55,above left] {$\mathbf r(t_*)$};
\node[orbit point] at (P) {};
\node[anchor=west] at ($(P)+(2,0)$) {$t_*$};
\end{tikzpicture}
\par
{\usebeamerfont{caption}Голубой эллипс --- оскулирующая орбита; штриховая дуга --- возмущённая траектория, схематически.\par}
\end{column}
\end{columns}
\end{frame}

% 8
\begin{frame}{Средние элементы}
\begin{block}{Средние элементы}
\justifying
Элементы, получаемые усреднением оскулирующих элементов по выбранному интервалу времени или угла с исключением короткопериодических вариаций. Они описывают долговременное поведение орбиты.
\end{block}
\begin{itemize}
\item \textbf{Интервал усреднения:} диапазон времени или выбранного угла.
\item \textbf{Переменная усреднения:} время, истинная, эксцентрическая или средняя аномалия, долгота.
\end{itemize}
\justifying
Произвольное сглаживание временного ряда не заменяет преобразования, согласованного с уравнениями эволюции средних элементов и периодическими поправками.
\end{frame}

% 9
\begin{frame}{Пример}
\setlength{\abovedisplayskip}{5pt}\setlength{\belowdisplayskip}{5pt}
\justifying
Пусть $x$ --- один из орбитальных элементов. Он может быть представлен в виде обрезанного ряда:
\[
x=x_0+\dot x_1(t-t_0)+K_1\cos(2\omega)
+K_2\sin(2\nu+\omega)+K_3\cos(2\nu).
\]
\par\noindent
{\renewcommand{\arraystretch}{1.2}
\begin{tabularx}{\dimexpr\textwidth-2pt\relax}{@{}p{.30\textwidth}>{\raggedright\arraybackslash}X@{}}
$x_0$ & Начальное значение.\\
$\dot x_1(t-t_0)$ & Вековое изменение со скоростью $\dot x_1$.\\
$K_1\cos(2\omega)$ & Долгопериодический член: медленный угол $\omega$.\\
$K_2\sin(2\nu+\omega)$ & Смешанный член: быстрая $\nu$ и медленная $\omega$.\\
$K_3\cos(2\nu)$ & Короткопериодический член: быстрая переменная $\nu$.
\end{tabularx}\par}
\medskip\justifying
$K_1$, $K_2$, $K_3$ --- постоянные, специфические для конкретного метода; обычно выражаются в виде сумм или произведений полиномов от $a$, $e$ или $i$.
\end{frame}

% 10
\begin{frame}{Однократное усреднение}
\setlength{\abovedisplayskip}{4pt}\setlength{\belowdisplayskip}{4pt}
\justifying
В приведённом примере усредним зависимость по быстрой фазе $\nu$ на $[0,2\pi]$. Параметры $t$, $\omega$, $x_0$, $\dot x_1$ и $K_i$ фиксируются.
\[
\bar x^{(1)}(t,\omega)=\frac{1}{2\pi}\int_0^{2\pi}x(t,\nu,\omega)\,d\nu.
\]
Короткопериодический и смешанный члены имеют нулевое среднее:
\par
{\centering $\int_0^{2\pi}\cos(2\nu)\,d\nu=0,\qquad
\int_0^{2\pi}\sin(2\nu+\omega)\,d\nu=0.$\par}
\begin{block}{Однократно усреднённый элемент}
\centering
$\displaystyle\bar x^{(1)}=x_0+\dot x_1(t-t_0)+K_1\cos(2\omega).$\par
\end{block}
\justifying
Сохраняются вековая и долгопериодическая части. Это усреднение по фазе: усреднение по времени на эллиптической орбите имеет другой вес, поскольку $\nu$ изменяется неравномерно.
\end{frame}

% 11
\begin{frame}{Двукратное усреднение}
\setlength{\abovedisplayskip}{4pt}\setlength{\belowdisplayskip}{4pt}
\justifying
Дополнительно усредним результат по медленному углу $\omega$ на $[0,2\pi]$, фиксируя $t$, $x_0$, $\dot x_1$ и $K_1$:
\[
\bar x^{(2)}(t)=\frac{1}{2\pi}\int_0^{2\pi}\bar x^{(1)}(t,\omega)\,d\omega.
\]
Долгопериодический член исчезает, поскольку
$\int_0^{2\pi}\cos(2\omega)\,d\omega=0$.
\begin{block}{Двукратно усреднённый элемент}
\centering
$\displaystyle\bar x^{(2)}=x_0+\dot x_1(t-t_0).$\par
\end{block}
\justifying
Сохраняется вековой член; два этапа усреднения --- по $\nu$ и $\omega$.\par\smallskip
Одному физическому состоянию КА могут соответствовать различные наборы средних элементов в разных теориях. Сопоставлять такие наборы следует через восстановленные оскулирующие элементы или декартово состояние.
\end{frame}

% 12
\begin{frame}{Метод возмущений: определение}
\setlength{\abovedisplayskip}{4pt}\setlength{\belowdisplayskip}{4pt}
\begin{block}{Определение}
\justifying
Метод возмущений --- совокупность математических приёмов построения приближённого решения задачи с малыми возмущающими воздействиями на основе известного решения соответствующей невозмущённой задачи.
\end{block}
\[
\dot{\mathbf x}=\mathbf f_0(\mathbf x,t)
+\varepsilon_1\mathbf f_1(\mathbf x,t)+\varepsilon_2\mathbf f_2(\mathbf x,t),
\qquad |\varepsilon_1|,|\varepsilon_2|\ll1.
\]
\justifying
$\mathbf f_0$ описывает невозмущённое движение; $\varepsilon_1\mathbf f_1$ и $\varepsilon_2\mathbf f_2$ --- возмущающие воздействия. Параметры $\varepsilon_1$, $\varepsilon_2$ безразмерны.\par\smallskip
При $\varepsilon_1=\varepsilon_2=0$ решение известно. Влияние возмущений определяется последовательным вычислением поправок к этому решению.
\end{frame}

% 13
\begin{frame}{Разложение решения и порядок теории}
Приближение к истинному решению может быть выражено как степенной ряд малых параметров:
\setlength{\abovedisplayskip}{4pt}\setlength{\belowdisplayskip}{4pt}
\[
\begin{aligned}
\mathbf x(t;\varepsilon_1,\varepsilon_2) =
\mathbf x_0(t) &+ \varepsilon_1\boldsymbol\alpha_1(t) + \frac{\varepsilon_1^2}{2}\boldsymbol\alpha_2(t) + \dots \\
&+ \varepsilon_2\boldsymbol\beta_1(t) + \frac{\varepsilon_2^2}{2}\boldsymbol\beta_2(t) + \dots \\
&+\varepsilon_1\varepsilon_2\boldsymbol\gamma_{11}(t) + \dots
\end{aligned}
\]
\justifying
$\mathbf x_0(t)$ --- невозмущённое решение; коэффициенты --- поправки, определяемые подстановкой ряда в уравнения движения и сопоставлением одинаковых порядков малости.

При одинаковом порядке $\varepsilon_1$, $\varepsilon_2$ первый порядок включает линейные члены; второй --- также квадратичные и смешанный член $\varepsilon_1\varepsilon_2$.

\end{frame}

% 14
\begin{frame}{Малые параметры}
\begin{block}{Примеры малых параметров}
\begin{itemize}
\item Коэффициенты гармоник гравитационного потенциала $J_2$, $J_3$, $C_{lm}$, $S_{lm}$.
\item Отношение больших полуосей КА и третьего тела $a/a_3$ для возмущений от третьих тел.
\item Базразмерная плотность для атмосферного сопротивления.
\item Безразмерное ускорение тяги ДУ для манёвров.
\item Безразмерный солнечный поток для давления солнечного излучения.
\end{itemize}
\end{block}

\end{frame}

% 15
\begin{frame}{Метод вариации параметра}
\setlength{\abovedisplayskip}{4pt}\setlength{\belowdisplayskip}{4pt}
\begin{block}{Определение}
\justifying
Метод вариации параметра представляет решение возмущённой задачи через известное решение невозмущённой задачи, в котором постоянные интегрирования заменены функциями времени.
\end{block}
\[
\ddot{\mathbf r}=-\frac{\mu}{r^3}\mathbf r,
\qquad \mathbf r=\mathbf R(\mathbf x,t),
\qquad \mathbf V(\mathbf x,t)=\frac{\partial\mathbf R}{\partial t}.
\]
\justifying
$\mathbf x=(a,e,i,\Omega,\omega,M_0)$ --- постоянные кеплерова решения; $M_0$ --- средняя аномалия на фиксированную эпоху.\par\smallskip
В возмущённой задаче полагают $\mathbf x=\mathbf x(t)$. Законы изменения параметров определяют из уравнений движения и дополнительных связей.
\end{frame}

% 16
\begin{frame}{Условие оскуляции}
\setlength{\abovedisplayskip}{4pt}\setlength{\belowdisplayskip}{4pt}
При $\mathbf r(t)=\mathbf R(\mathbf x(t),t)$ полная производная равна
\[
\dot{\mathbf r}=\frac{\partial\mathbf R}{\partial t}
+\sum\nolimits_{j=1}^{6}\frac{\partial\mathbf R}{\partial x_j}\dot x_j.
\]
\begin{block}{Условие оскуляции}
\centering
$\displaystyle\sum\nolimits_{j=1}^{6}\frac{\partial\mathbf R}{\partial x_j}\dot x_j=\mathbf0.$
\par\smallskip\justifying
Кеплерово представление воспроизводит фактические положение и скорость КА в каждый момент времени: $\dot{\mathbf r}=\mathbf V(\mathbf x(t),t)$.
\end{block}
\justifying
Три координаты не определяют эволюцию шести параметров однозначно. Условие оскуляции задаёт три дополнительные связи.
\end{frame}

% 17
\begin{frame}{Уравнения изменения элементов}
\setlength{\abovedisplayskip}{4pt}\setlength{\belowdisplayskip}{4pt}
\[
\ddot{\mathbf r}=-\frac{\mu}{r^3}\mathbf r+\mathbf a_{\text{возм}},
\qquad
\begin{bmatrix}D_{\mathbf x}\mathbf R\\D_{\mathbf x}\mathbf V\end{bmatrix}
\dot{\mathbf x}=\begin{bmatrix}\mathbf0\\\mathbf a_{\text{возм}}\end{bmatrix}.
\]
\justifying
$D_{\mathbf x}\mathbf R$ и $D_{\mathbf x}\mathbf V$ --- производные положения и скорости по параметрам. Верхние три уравнения задают оскуляцию, нижние три --- действие возмущающего ускорения.
\begin{block}{Система для шести параметров}
\centering
$\dot{\mathbf x}=\mathbf G(\mathbf x,t)\,\mathbf a_{\text{возм}}.$
\par\smallskip\justifying
Матрица преобразования должна быть невырожденной. При замене $M_0$ текущей аномалией $M$ в её уравнение входит кеплеровский член $n$.
\end{block}
\justifying
Полученную систему решают аналитически или численно.
\end{frame}

% 18
\begin{frame}{Формы Лагранжа и Гаусса}
\setlength{\abovedisplayskip}{4pt}\setlength{\belowdisplayskip}{4pt}
\begin{columns}[T,onlytextwidth]
\begin{column}{.48\textwidth}
\begin{block}{Форма Лагранжа}
\justifying
Потенциальное возмущение задаётся возмущающей функцией $\mathcal R(\mathbf r,t)$:
\[
\mathbf a_{\text{возм}}=\nabla_{\mathbf r}\mathcal R.
\]
Производные $\mathcal R$ по элементам определяют их скорости:
\[
\dot\Omega=\frac{1}{na^2\sqrt{1-e^2}\sin i}\,
\frac{\partial\mathcal R}{\partial i}.
\]
\end{block}
\end{column}
\begin{column}{.48\textwidth}
\begin{block}{Форма Гаусса}
\justifying
Возмущение задаётся компонентами ускорения:
\[
\mathbf a_{\text{возм}}=a_R\mathbf e_R+a_T\mathbf e_T+a_N\mathbf e_N.
\]
Применима к потенциальным и непотенциальным силам. Например,
\[
\dot\Omega=\frac{r\sin u}{h\sin i}\,a_N.
\]
\end{block}
\end{column}
\end{columns}
\smallskip
$u=\omega+\nu$, $h=\lvert\mathbf r\times\mathbf v\rvert$, $n=\sqrt{\mu/a^3}$.
\par\smallskip
Направления $R,T,N$: радиальное, трансверсальное и нормальное к орбите.
\end{frame}

% 19
\begin{frame}{Возмущение $J_2$ в уравнениях движения}
\setlength{\abovedisplayskip}{4pt}\setlength{\belowdisplayskip}{4pt}
\justifying
Вторая зональная гармоника $J_2$ характеризует полярное сжатие Земли;
$R_\oplus$ --- её экваториальный радиус.
\[
\mathcal R_{J_2}=\frac{\mu J_2R_\oplus^2}{2r^3}
\bigl(1-3\sin^2i\sin^2u\bigr),
\qquad u=\omega+\nu,\quad r=\frac{a(1-e^2)}{1+e\cos\nu}.
\]
\begin{block}{Усреднение по средней аномалии}
\centering
$\displaystyle
\overline{\mathcal R}_{J_2}
=\frac{1}{2\pi}\int_0^{2\pi}\mathcal R_{J_2}\,dM
=\frac{\mu J_2R_\oplus^2}{4a^3(1-e^2)^{3/2}}
\bigl(3\cos^2i-1\bigr).$
\par\smallskip\justifying
При усреднении $a,e,i,\Omega,\omega$ считаются постоянными;
радиус $r$ и истинная аномалия $\nu$ изменяются вдоль кеплеровой орбиты.
\end{block}
\justifying
Из уравнений Лагранжа с $\overline{\mathcal R}_{J_2}$ получают вековые скорости элементов.
\end{frame}

% 20
\begin{frame}{Вековые изменения вследствие $J_2$}
\setlength{\abovedisplayskip}{3pt}\setlength{\belowdisplayskip}{3pt}
Первое приближение по $J_2$; элементы в формулах --- средние.
\begin{columns}[T,onlytextwidth]
\begin{column}{.62\textwidth}
\begin{block}{Вековые скорости}
\centering
$\displaystyle\dot\Omega_{J_2}
=-\frac32J_2n\bigl(R_\oplus/p\bigr)^2\cos i,$
\par\smallskip
$\displaystyle\dot\omega_{J_2}
=\frac34J_2n\bigl(R_\oplus/p\bigr)^2(5\cos^2i-1),$
\par\smallskip
$\displaystyle\dot M_{J_2}
=\frac34J_2n\bigl(R_\oplus/p\bigr)^2
\sqrt{1-e^2}\,(3\cos^2i-1).$
\par\smallskip
$\displaystyle\dot M=n+\dot M_{J_2}.$
\end{block}
$n=\sqrt{\mu/a^3}$,\quad $p=a(1-e^2)$.
\par\smallskip
При $n$ в с$^{-1}$ скорости выражены в рад/с.
\end{column}
\begin{column}{.34\textwidth}
\textbf{Прецессия плоскости}
\par\smallskip
Для прямых орбит $\dot\Omega_{J_2}<0$, для обратных --- $\dot\Omega_{J_2}>0$.
При $i=90^\circ$ имеем $\dot\Omega_{J_2}=0$.
\par\medskip
\textbf{Вращение линии апсид}
\par\smallskip
$\dot\omega_{J_2}=0$ при критических наклонениях
$i\approx63{,}4^\circ$ и $116{,}6^\circ$.
\par\medskip
Вековые скорости $a,e,i$:
$\dot a=\dot e=\dot i=0$.
\end{column}
\end{columns}
\end{frame}

% 21
\begin{frame}{Основное влияние других возмущений}
\renewcommand{\arraystretch}{1.10}
\par\noindent
\begin{tabularx}{\dimexpr\textwidth-2pt\relax}{@{}>{\raggedright\arraybackslash}p{32mm}>{\raggedright\arraybackslash}X@{}}
\toprule
\textbf{Возмущение} & \textbf{Изменения орбитальных элементов}\\
\midrule
Сопротивление атмосферы &
Уменьшение $a$ и периода обращения; увеличение $n$.
При преимущественном торможении вблизи перицентра уменьшаются
высота апоцентра и эксцентриситет.\\
\midrule
Притяжение Луны и Солнца &
Прецессия плоскости орбиты и линии апсид; долгопериодические
изменения $e,i,\Omega,\omega$.
В нерезонансном двукратно усреднённом приближении $a$ постоянно.\\
\midrule
Давление солнечного излучения &
Периодические изменения элементов, прежде всего $e$ и $\omega$.
Вблизи резонансов возможны значительные долгопериодические изменения.
Величина эффектов зависит от отношения площади к массе, ориентации КА и затенения.\\
\bottomrule
\end{tabularx}
\end{frame}

% 22
\begin{frame}{Точность прогноза и измерений}
\begin{columns}[T,onlytextwidth]
\begin{column}{.65\textwidth}\centering
% Перерисовка исторической схемы Vallado по исходному рисунку лекции 11.
% Положения кривых восстановлены графически, обе шкалы логарифмические.
\begin{tikzpicture}[course diagram,x=1mm,y=.95mm]
\path[use as bounding box] (-11,-7) rectangle (85,53);
\node[anchor=south] at (36,51) {Погрешность положения};
\foreach \y/\lab in {0/{10 см},10/{1 м},20/{10 м},30/{100 м},40/{1 км},50/{10 км}}{
 \draw[aux line] (0,\y)--(75,\y);
 \node[anchor=east,inner sep=1mm] at (0,\y) {\lab};
}
\foreach \x/\lab in {0/{100},25/{1000},50/{10\,000},75/{100\,000}}{
 \draw[aux line] (\x,0)--(\x,50);
 \node[below,inner sep=1mm] at (\x,0) {$\lab$};
}
% Промежуточные деления логарифмических шкал.
\foreach \d in {0,1,2}{\foreach \k in {2,...,9}{
 \pgfmathsetmacro{\tickx}{25*(\d+ln(\k)/ln(10))}
 \draw[coursemuted,line width=.35pt] (\tickx,0)--(\tickx,-.7);
}}
\foreach \d in {0,...,4}{\foreach \k in {2,...,9}{
 \pgfmathsetmacro{\ticky}{10*(\d+ln(\k)/ln(10))}
 \draw[coursemuted,line width=.35pt] (0,\ticky)--(-.6,\ticky);
}}
\draw[courseink] (0,50)--(0,0)--(75,0);
\node[anchor=north,inner sep=0pt] at (37.5,-4.5) {$H$, км};
% 1--3: аналитический прогноз, аппроксимация наблюдений, теория.
\draw[courseaccent,line width=1.1pt]
 (3.7,50) .. controls (5.5,45) and (18,43.7) .. (32.5,43.7)
 .. controls (45,43.7) and (58,45.2) .. (64,49.8);
\draw[courseaccent,line width=.65pt]
 (3.8,45.2) .. controls (8,41.6) and (20,40) .. (35,40)
 .. controls (50,40) and (63,41.5) .. (68.4,45.4);
\draw[courseaccent,line width=1.3pt] (2.5,39)--(71,39);
% 4--5: РЛС низкой точности и интерферометр.
\draw[coursemuted,dashed,line width=.65pt]
 (3,34.9) .. controls (34,34.5) and (50,34.4) .. (72.5,36.5);
\draw[coursemuted,dashed,line width=.65pt] (2.5,33)--(45,33);
% 6--9: численный прогноз, РЛС высокой точности, теория, лазерные измерения.
\draw[courseink,line width=1.1pt]
 (4.5,21.1) .. controls (10,19.6) and (20,19.4) .. (32,19.4)
 .. controls (47,19.4) and (62,21.1) .. (72.5,23);
\draw[coursemuted,dashed,line width=.65pt]
 (3.4,19.1) .. controls (31,18.4) and (49,18.8) .. (73,21.3);
\draw[courseink,line width=1.3pt] (3,16)--(70,16);
\draw[coursemuted,dashed,line width=.65pt] (3,7.8)--(67.5,7.8);
% Номера размещены с фиксированными промежутками; тонкие линии --- выноски.
\tikzset{accuracy leader/.style={draw=coursemuted,line width=.4pt}}
\draw[accuracy leader] (32.5,43.7)--(32.5,45.4);
\node[anchor=south,inner sep=.3mm,text=courseaccent] at (32.5,45.4) {1};
\draw[accuracy leader] (68.4,45.4)--(72,47);
\node[anchor=south,inner sep=.3mm,text=courseaccent] at (72,47) {2};
\draw[accuracy leader] (45,33)--(48,31);
\node[anchor=west,inner sep=.6mm] at (48,31) {5};
\foreach \xe/\ye/\yl/\num in {71/39/40/3,72.5/36.5/35/4,72.5/23/25/6,73/21.3/20/7,70/16/15/8,67.5/7.8/7.8/9}{
 \draw[accuracy leader] (\xe,\ye)--(79,\yl)--(80,\yl);
 \node[anchor=west,inner sep=.6mm] at (80,\yl) {\num};
}
\end{tikzpicture}

\end{column}
\begin{column}{.33\textwidth}
{\usebeamerfont{caption}\raggedright
\textcolor{courseaccent}{1. Аналитический прогноз}\par
\textcolor{courseaccent}{2. Аналитическая аппроксимация наблюдений}\par
\textcolor{courseaccent}{3. Аналитическая теория}\par
4. РЛС низкой точности\par
5. Интерферометр\par
6. Численный прогноз\par
7. РЛС высокой точности\par
8. Численная теория\par
9. Лазерные измерения\par}
\end{column}
\end{columns}
\end{frame}

% 23
\begin{frame}{Численный прогноз движения}
\setlength{\abovedisplayskip}{4pt}\setlength{\belowdisplayskip}{4pt}
\justifying
\textbf{Численный прогноз} состоит в приближённом решении задачи Коши для уравнений движения последовательным вычислением состояния на шагах интегрирования.\par
\[
\dot{\mathbf r}=\mathbf v,\qquad
\dot{\mathbf v}=-\frac{\mu}{r^3}\mathbf r+\mathbf a_{\text{возм}}(\mathbf r,\mathbf v,t),
\qquad (\mathbf r,\mathbf v)\big|_{t_0}=(\mathbf r_0,\mathbf v_0).
\]
\begin{itemize}
\item Начальные данные: положение, скорость и эпоха; параметры моделей сил задаются отдельно.
\item На каждом шаге вычисляются возмущающие ускорения и приращение состояния.
\item Результат: положение и скорость на заданные эпохи; орбитальные элементы получают преобразованием состояния.
\end{itemize}
\end{frame}

% 24
\begin{frame}{Численные методы}
\begin{columns}[T,onlytextwidth]
\begin{column}{.48\textwidth}
\begin{block}{Возможности}
\begin{itemize}\setlength{\itemsep}{4pt}
\item Сложные модели сил и их явная зависимость от времени, положения и скорости.
\item Контроль численной погрешности выбором метода, шага и допусков.
\item Учёт изменений режима движения, включая манёвры, на заданных эпохах.
\end{itemize}
\end{block}
\end{column}
\begin{column}{.48\textwidth}
\begin{block}{Ограничения}
\begin{itemize}\setlength{\itemsep}{4pt}
\item Шаг определяется характером изменения правых частей, включая короткую периодику.
\item Затраты возрастают с длительностью прогноза и сложностью моделей сил.
\item Уменьшение численной погрешности не устраняет ошибки начальных данных и моделей сил.
\end{itemize}
\end{block}
\end{column}
\end{columns}
\end{frame}

% 25
\begin{frame}{Построение аналитического решения}
\setlength{\abovedisplayskip}{3pt}\setlength{\belowdisplayskip}{3pt}
\justifying
\textbf{Аналитический прогноз} использует приближённое решение, явно выраженное через время и начальные параметры. В методе возмущений решение строится по порядкам малого параметра.\par
\[
\dot{\mathbf x}=\mathbf f_0(\mathbf x,t)+\varepsilon\mathbf f_1(\mathbf x,t),
\qquad \mathbf x=\mathbf x_0+\varepsilon\mathbf x_1+O(\varepsilon^2).
\]
\textbf{Уравнения нулевого и первого порядков}\par
\[
\dot{\mathbf x}_0=\mathbf f_0(\mathbf x_0,t),\qquad
\dot{\mathbf x}_1=D_{\mathbf x}\mathbf f_0(\mathbf x_0,t)\,\mathbf x_1
+\mathbf f_1(\mathbf x_0,t).
\]
\begin{enumerate}\setlength{\itemsep}{2pt}
\item Подставить невозмущённое решение в уравнения поправок.
\item Проинтегрировать сохраняемые вековые и периодические члены.
\item Определить постоянные из начальных данных и вычислить состояние.
\end{enumerate}
\end{frame}

% 26
\begin{frame}{Аналитический прогноз с учётом $J_2$}
\setlength{\abovedisplayskip}{3pt}\setlength{\belowdisplayskip}{3pt}
\justifying
Простейшее приближение сохраняет только вековые эффекты первого порядка по $J_2$. Исходные средние элементы заданы на эпоху $t_0$; $\Delta t=t-t_0$.
\par\smallskip
\textbf{Изменение средних элементов}\par
\[
\begin{gathered}
\bar a=\bar a_0,\qquad \bar e=\bar e_0,\qquad \bar i=\bar i_0,\\
\bar\Omega=\bar\Omega_0+\dot\Omega_{J_2}\Delta t,\qquad
\bar\omega=\bar\omega_0+\dot\omega_{J_2}\Delta t,\\
\bar M=\bar M_0+(n+\dot M_{J_2})\Delta t.
\end{gathered}
\]
\justifying
Скорости вычисляются по $\bar a_0,\bar e_0,\bar i_0$. Из $\bar M$ находят $E$ по уравнению Кеплера, затем $\nu$ и декартово состояние.\par\smallskip
Короткопериодические поправки не восстановлены. Уменьшение шага выдачи результатов не устраняет вызванную этим погрешность состояния.\par
\end{frame}

% 27
\begin{frame}{Аналитические методы}
\begin{columns}[T,onlytextwidth]
\begin{column}{.48\textwidth}
\begin{block}{Возможности}
\begin{itemize}\setlength{\itemsep}{4pt}
\item Вычисление состояния на заданную эпоху без пошагового интегрирования.
\item Явная зависимость вековых скоростей и периодических поправок от элементов орбиты.
\item Повторное применение выражений в установленной области параметров.
\end{itemize}
\end{block}
\end{column}
\begin{column}{.48\textwidth}
\begin{block}{Ограничения}
\begin{itemize}\setlength{\itemsep}{4pt}
\item Состав сил и допустимые орбиты заданы конкретной теорией.
\item Учёт новых сил и повышение порядка требуют вывода дополнительных выражений.
\item Ошибки вековых членов накапливаются; исходные средние элементы должны соответствовать теории.
\end{itemize}
\end{block}
\end{column}
\end{columns}
\end{frame}

% 28
\begin{frame}{Теория Козая}
\justifying
Аналитическая теория движения спутника в поле несферической Земли с зональными гармониками $J_2$, $J_3$, $J_4$ (1959 г.). Сопротивление атмосферы и притяжение третьих тел в исходной теории отсутствуют.\par
\begin{enumerate}\setlength{\itemsep}{3pt}
\item Уравнения Лагранжа интегрируются с разделением вековых, долго- и короткопериодических составляющих.
\item Периодические поправки определяются разложением по возмущающим коэффициентам; в вековых скоростях учитываются также члены $J_2^2$.
\item На требуемую эпоху средняя эволюция дополняется периодическими поправками, после чего восстанавливается состояние КА.
\end{enumerate}
\justifying
Постоянные интегрирования и средние элементы определяются принятым способом разделения составляющих решения.
\end{frame}

% 29
\begin{frame}{Теория Брауэра}
\justifying
В теории Брауэра (1959 г.) два канонических преобразования переменных Делоне исключают из гамильтониана $M$ и $\omega$. Решение содержит вековую эволюцию и периодические поправки, восстанавливаемые обратными преобразованиями.\par\smallskip
\textbf{\color{courseaccent}Модель возмущений}
\begin{itemize}\setlength{\itemsep}{2pt}
\item Зональные гармоники гравитационного поля Земли $J_2$, $J_3$, $J_4$, $J_5$.
\item Для $J_2$: вековые изменения --- до второго порядка ($J_2^2$); коротко- и долгопериодические поправки --- первого порядка.
\item Сопротивление атмосферы, притяжение Луны и Солнца, давление солнечного излучения и неосесимметричные гармоники не учитываются.
\end{itemize}
\justifying
При малых $e$, $i$ и близости к критическим наклонениям требуются модификации. Средние элементы Брауэра и Козая различаются; сравнивают восстановленные положения и скорости.
\end{frame}

% 30
\begin{frame}{SGP4 и исходные элементы TLE}
\justifying
\textbf{SGP4} (Simplified General Perturbations 4) --- модель прогноза, согласованная с двухстрочными наборами орбитальных элементов TLE.\par\smallskip
\textbf{\color{courseaccent}Модель возмущений}
\begin{itemize}\setlength{\itemsep}{2pt}
\item Зональные гармоники $J_2$, $J_3$, $J_4$; в вековых скоростях сохраняются члены $J_2^2$.
\item Упрощённое сопротивление атмосферы с параметром $B^*$ из TLE.
\item При периоде $T\geq225$ мин: притяжение Луны и Солнца и резонансные члены неосесимметричного геопотенциала (резонансы $1{:}1$ и $2{:}1$ с вращением Земли).
\end{itemize}
\justifying
Давление солнечного излучения в стандартной SGP4 явно не учитывается.\par\smallskip
\textbf{Исходные данные:} средние элементы TLE на заданную эпоху и $B^*$. \textbf{Результат:} положение и скорость в системе TEME.\par\smallskip
Элементы TLE нельзя непосредственно подставлять в другую теорию движения.\par
\end{frame}

% 31
\begin{frame}{Принцип полуаналитического прогноза}
\begin{block}{Определение}
\justifying
Полуаналитический прогноз основан на численном интегрировании уравнений для средних элементов и восстановлении оскулирующих элементов с помощью периодических поправок, выделенных методами теории возмущений.
\end{block}
\begin{enumerate}
\item Усреднением уравнений движения выделяют вековые и долгопериодические изменения элементов.
\item Полученную систему интегрируют численно. Шаг определяется изменением усреднённых правых частей.
\item В заданные моменты времени вычисляют короткопериодические поправки и восстанавливают положение и скорость КА.
\end{enumerate}
\end{frame}

% 32
\begin{frame}{Уравнения для средних элементов}
\setlength{\abovedisplayskip}{5pt}
\setlength{\belowdisplayskip}{5pt}
\setlength{\abovedisplayshortskip}{0pt}
\setlength{\belowdisplayshortskip}{4pt}
\justifying
Пусть $\widehat{\mathbf y}$ --- пять медленных оскулирующих элементов, $\hat\lambda$ --- быстрая средняя долгота. Рассмотрим систему
\[
\dot{\widehat{\mathbf y}}=\varepsilon\mathbf F(\widehat{\mathbf y},\hat\lambda),
\qquad
\dot{\hat\lambda}=n(\widehat{\mathbf y})+O(\varepsilon).
\]
\medskip
\textbf{\color{courseaccent}Первое приближение}\par
\justifying
При усреднении медленные элементы фиксируют. Если $\mathbf F$ периодична по $\lambda$ с периодом $2\pi$, то
\[
\langle\mathbf F\rangle(\bar{\mathbf y})
=\frac{1}{2\pi}\int_0^{2\pi}\mathbf F(\bar{\mathbf y},\lambda)\,d\lambda,
\qquad
\dot{\bar{\mathbf y}}=\varepsilon\langle\mathbf F\rangle(\bar{\mathbf y})
+O(\varepsilon^2).
\]
\justifying
Усредняется правая часть уравнений. Эволюцию средних элементов определяют численным решением полученной системы.
\end{frame}

% 33
\begin{frame}{Короткопериодические поправки}
\setlength{\abovedisplayskip}{5pt}
\setlength{\belowdisplayskip}{5pt}
\setlength{\abovedisplayshortskip}{0pt}
\setlength{\belowdisplayshortskip}{4pt}
\justifying
Связь медленных оскулирующих и средних элементов в первом порядке задаётся преобразованием
\[
\widehat{\mathbf y}=\bar{\mathbf y}
+\varepsilon\boldsymbol\eta(\bar{\mathbf y},\bar\lambda)
+O(\varepsilon^2).
\]
Подстановка в исходные уравнения и сравнение членов порядка $\varepsilon$ дают
\[
n(\bar{\mathbf y})\frac{\partial\boldsymbol\eta}{\partial\bar\lambda}
=\mathbf F(\bar{\mathbf y},\bar\lambda)
-\langle\mathbf F\rangle(\bar{\mathbf y}),
\qquad
\langle\boldsymbol\eta\rangle=\mathbf0.
\]
\begin{block}{Определение поправки}
\justifying
Поправку находят интегрированием $(\mathbf F-\langle\mathbf F\rangle)/n$ по быстрой переменной. Постоянную интегрирования определяют из условия нулевого среднего значения поправки.
\end{block}
\end{frame}

% 34
\begin{frame}{DSST: переменные и уравнения}
\setlength{\abovedisplayskip}{5pt}
\setlength{\belowdisplayskip}{5pt}
\setlength{\abovedisplayshortskip}{0pt}
\setlength{\belowdisplayshortskip}{4pt}
\textbf{Draper Semianalytical Satellite Theory}\par
Средние равноденственные элементы для наклонений $i\ne180^\circ$:
\[
\begin{gathered}
\mathbf x=(a,h,k,p,q,\lambda)^{\mathsf T},\qquad
h=e\sin(\omega+\Omega),\quad k=e\cos(\omega+\Omega),\\
p=\tan\frac{i}{2}\sin\Omega,\quad q=\tan\frac{i}{2}\cos\Omega,
\qquad\lambda=M+\omega+\Omega.
\end{gathered}
\]
Уравнения средних элементов имеют структуру
\[
\dot{\bar x}_i=n(\bar a)\delta_{i6}
+\sum_{j=1}^{m}\varepsilon^j A_{ij}(\bar{\mathbf x},t),
\qquad i=1,\ldots,6.
\]
\justifying
$A_{ij}$ --- усреднённые скорости порядка $j$; $m$ --- порядок теории;
$\delta_{i6}$ --- символ Кронекера.
Средняя долгота сохраняет быстрое изменение $n(\bar a)$ и определяет орбитальную фазу.
\end{frame}

% 35
\begin{frame}{DSST: начальные средние элементы}
\setlength{\abovedisplayskip}{5pt}
\setlength{\belowdisplayskip}{5pt}
\setlength{\abovedisplayshortskip}{0pt}
\setlength{\belowdisplayshortskip}{4pt}
\justifying
Оскулирующие элементы $\hat{\mathbf x}_0$ на эпоху $t_0$ связаны со средними элементами неявным соотношением
\[
\hat{\mathbf x}_0=\bar{\mathbf x}_0
+\Delta\mathbf x_{\mathrm{SP}}(\bar{\mathbf x}_0,t_0).
\]
\medskip
\textbf{\color{courseaccent}Итерационное преобразование}\par
\[
\bar{\mathbf x}_0^{(0)}=\hat{\mathbf x}_0,
\qquad
\bar{\mathbf x}_0^{(k+1)}=\hat{\mathbf x}_0
-\Delta\mathbf x_{\mathrm{SP}}(\bar{\mathbf x}_0^{(k)},t_0).
\]
\justifying
На каждой итерации поправки пересчитывают по текущему приближению средних элементов. Сходимость контролируют по восстановлению заданного оскулирующего состояния.\par
\medskip
\justifying
$\Delta\mathbf x_{\mathrm{SP}}$ --- короткопериодические поправки. Их модели сил и порядок приближения согласуют с усреднёнными уравнениями.
\end{frame}

% 36
\begin{frame}{DSST: последовательность прогноза}
\centering
\begin{tikzpicture}[course diagram,x=1mm,y=1mm,
 stage/.style={draw=courserule,fill=coursetint,rounded corners=1.2mm,
 text width=64mm,minimum height=12mm,align=center,inner sep=2mm,
 font=\fontsize{11}{14}\selectfont},
 arrow/.style={->,draw=courseaccent,line width=.9pt}]
\node[stage] (initial) at (0,0) {Начальное состояние\\$\mathbf r_0,\ \mathbf v_0,\ t_0$};
\node[stage] (mean) at (76,0) {Преобразование\\$\hat{\mathbf x}_0\longrightarrow\bar{\mathbf x}_0$};
\node[stage] (integrate) at (76,-21) {Интегрирование уравнений\\средних элементов};
\node[stage] (periodic) at (0,-21) {Восстановление элементов\\$\hat{\mathbf x}(t)=\bar{\mathbf x}(t)+\Delta\mathbf x_{\mathrm{SP}}(t)$};
\node[stage] (state) at (0,-42) {Прогнозируемое состояние\\$\mathbf r(t),\ \mathbf v(t)$};
\draw[arrow] (initial.east)--(mean.west);
\draw[arrow] (mean.south)--(integrate.north);
\draw[arrow] (integrate.west)--node[above,text=coursemuted,font=\fontsize{9.5}{11.5}\selectfont] {$\bar{\mathbf x}(t)$} (periodic.east);
\draw[arrow] (periodic.south)--(state.north);
\node[align=left,text width=63mm,anchor=north west,
 font=\fontsize{11}{14}\selectfont] at (42,-35)
 {Поправки вычисляются на требуемые эпохи с учётом орбитальной фазы.};
\end{tikzpicture}
\end{frame}

% 37
\begin{frame}{Полуаналитические методы: сравнение}
\vspace{-1.5mm}
\begin{columns}[T,onlytextwidth]
\begin{column}{.48\textwidth}
\begin{block}{Возможности}
\begin{itemize}\setlength{\itemsep}{1.5pt}
\item Шаг интегрирования может превышать орбитальный период.
\item Снижение затрат при длительном прогнозе с редкой выдачей состояний.
\item Учёт короткой периодики без интегрирования движения на каждом витке.
\end{itemize}
\end{block}
\end{column}
\begin{column}{.48\textwidth}
\begin{block}{Ограничения}
\begin{itemize}\setlength{\itemsep}{1.5pt}
\item Погрешности усечения теории и численного решения.
\item Необходимость вывода средних скоростей и поправок для каждой модели сил.
\item Отдельный учёт резонансов и быстрых непериодических воздействий.
\end{itemize}
\end{block}
\end{column}
\end{columns}
\smallskip
\justifying
Импульс задаёт скачок оскулирующего состояния; после него преобразование к средним элементам выполняют повторно.
\end{frame}

% 38
\begin{frame}{Источники}
\setbeamertemplate{bibliography item}{\insertbiblabel}
\begin{thebibliography}{9}
\bibitem{Vallado}
Vallado D. A. Fundamentals of Astrodynamics and Applications. 2001.

\bibitem{Kozai}
Kozai Y. The Motion of a Close Earth Satellite. The Astronomical Journal. 1959. Vol. 64. P. 367--377. \href{https://doi.org/10.1086/107957}{DOI: 10.1086/107957}.

\bibitem{Brouwer}
Brouwer D. Solution of the Problem of Artificial Satellite Theory without Drag. The Astronomical Journal. 1959. Vol. 64. P. 378--397. \href{https://doi.org/10.1086/107958}{DOI: 10.1086/107958}.

\bibitem{Danielson}
Danielson D. A., Sagovac C. P., Neta B., Early L. W. – Monterey: Naval Postgraduate School, 1995.
\end{thebibliography}
\end{frame}

\end{document}
